\section{Integral Curves}

To prove the existence of integral curve, one needs to solve an ODE in the model space. \todo{expand on this}.
The typical way to do this is via the \emph{Picard-Lindelöf} theorem (also called the \emph{Cauchy-Lipschitz} theorem). 
This and other similar theorems about ODE are typically formulated for open sets. 
Since we work with vector fields defined on sets that need not be open, we need to adapt. 
In a lot of differential geometry textbooks, continuous differentiability and similar notions in not-open sets are defined to mean that there exist local extensions to open neighbourhoods that have this property. \todo{discuss this in prerequisites section}
However, \mathlib has chosen to not do this. \todo{cite Gouezels paper}
So, if we wanted to extend our functions, we would need to have a theorem that permits this. That theorem is the \emph{Whitney extension theorem}. 
Unfortunately, the proof of this theorem is quite long and thus did not appear to be in scope of this thesis. 
Additionally, it seems that \mathlib has found a lot of beautiful general statements by avoiding this theorem and instead develop a theory that also works for not-open sets.

So this is the route that we also chose to go for. 
Thus we also needed a version of the Picard-Lindelöf theorem for our use case. 
\todo{say that this is all based on the mathlib implementation}


We first collect some conditions for our theorem: 

\begin{defi}
Let $E$ be a normed commutative group under addition and a normed space over $\bR$. 
Let $s \subseteq E$ a set, $\minsub{t}, \maxsub{t} \in \bR$, $f \colon [\minsub{t}, \maxsub{t}] \times s \to E$ a function, $x_0 \in E$ and $a, r, L, K \in \bR_{\ge 0}$. 
Then we say that $f$ is \emph{Picard-Lindelöf} with respect to $x_0, a, r, L, K$ if the following conditions hold: 
\begin{enumerate}
    \item \label{PicLindI} \emph{Lipschitz condition}: for all $t \in [\minsub{t}, \maxsub{t}]$, the function $f_t \colon s \to E, x \mapsto f ((t, x))$ is Lipschitz on $s \cap \closedBall{a}{x_0}$ with Lipschitz constant $L$. 
    \item \label{PicLindII} \emph{Continuity}: for all $x \in s \cap \closedBall{a}{x_0}$, the function $f^x \colon [\minsub{t}, \maxsub{t}] \to E, t \mapsto f ((t, x))$ is continuous. 
    \item \label{PicLindIII} \emph{Norm condition}: for all $t \in [\minsub{t}, \maxsub{t}]$ and $x \in s \cap \closedBall{a}{x_0}$, $\norm{f ((t, x))} \le L$.
    \item \label{PicLindIV} \emph{Time condition}: We have $L (\maxsub{t} - \minsub{t}) \le a - r$. 
    \item \label{PicLindV} \emph{Set conditions}: The set $s$ is convex and closed. 
    \item \label{PicLindVI}\emph{Direction of vector field}: for all $v \in E$ if there is a $t \in [\minsub{t}, \maxsub{t}]$ such that $\norm{v - f ((t, x_0))} \le Ka$, then $v$ is inward pointing with respect to $s$ at $x_0$.
    \item \label{PicLindVII}\emph{Shape condition:} for all vectors $v \in E$ that are derivable within $s$ at $x_0$ and fulfil $\norm{v} \le a$, we get $x_0 + v \in s$. 
\end{enumerate}
\end{defi}

For the rest of the section let us fix a normed commutative additive group $E$ which is also a normed space over $\bR$ and a complete space. Additionally fix $\minsub{t} \le \maxsub{t}$ in $\bR$, $s \subseteq E$, a vector field $f \colon [\minsub{t}, \maxsub{t}] \times s \to E$, $x_0 \in s$ and $a, r, L, K \in \bR_{\ge}$. 

The theorem that we want to prove is the following: 

\begin{lem}
Assume that $f$ is Picard-Lindelöf with respect to $x_0, a, r, L, K$. 
The there is a differentiable function $\alpha \colon [\minsub{t}, \maxsub{t}] \to E$ such that for all $t \in [\minsub{t}, \maxsub{t}]$ it holds that $\alpha'(t) = f ((t, \alpha (t)))$ and $\alpha(t) \in s$. 
\end{lem}

Like in the proof of the usual Picard-Lindelöf theorem, we prove this by defining a contraction on a function space and then showing that the fixed point fulfils the conditions required in the theorem. 

So let us first define the function space: 

\begin{defi}
We define $\fM$ to be the space of functions $g \colon [\minsub{t}, \maxsub{t}] \to E$ such that
\begin{enumerate}
    \item $g$ is Lipschitz continuous with constant $L$,
    \item $g (\minsub{t}) \in \closedBall{r}{x_0}$, and 
    \item $g ([\minsub{t}, \maxsub{t}]) \subseteq s$.
\end{enumerate}
\end{defi}

Since we later want to apply the fact that every contraction in a complete space has a fixed point, we first need to show that this space is complete. 

\begin{lem}
If $s$ is closed, then $\fM$ is a complete space. 
\end{lem}
\begin{proof}
Since $\fM$ is a subspace of the complete space $\contMaps{[\minsub{t}, \maxsub{t}]}{E}$ it suffices to show that it is closed. 
We can express $\fM$ as the intersection of the following three sets: 
\begin{equation*}
    \begin{split}
    \{g \in \contMaps{[\minsub{t}, \maxsub{t}]}{E} &\mid g \text{ is Lipschitz with constant } L\} \\
    \{g \in \contMaps{[\minsub{t}, \maxsub{t}]}{E} &\mid \norm{g (\minsub{t}) - x_0} \le r\} \\ 
    \{g \in \contMaps{[\minsub{t}, \maxsub{t}]}{E} &\mid g ([\minsub{t}, \maxsub{t}]) \subseteq s\}
    \end{split}
\end{equation*}

In any metric space, the set of function that a Lipschitz with a specific constant is closed. Thus the first set is closed. 
The second set is closed since the function $g \mapsto \norm{g (\minsub{t}) - x_0}$ is continuous. 
And the last set is closed because $s$ is closed. 

Thus, as the intersection of these three sets, $\fM$ is also closed. 

\end{proof}

Let us now define the relevant contraction: 

\begin{defi}
We define the function $C \colon \fM \to ([\minsub{t}, \maxsub{t}] \to E)$ as 
\begin{equation*}
    \begin{split}
        C(g) \colon [\minsub{t}, \maxsub{t}] &\to E \\
        t &\mapsto x_0 + \int_{\minsub{t}}^t f(x, g(x)) \diff x
    \end{split}
\end{equation*}
\end{defi}

A priori, we do not know that this function again maps to $\fM$.
So let us prove this. 

We first need some helper lemmas: 

\begin{lem}\label{lem:funInBall}
Let $f \in \fM$ and $L (\maxsub{t} - \minsub{t}) \le a - r$. Then for all $t \in [\minsub{t}, \maxsub{t}]$, we get $f (t) \in \closedBall{a}{x_0}$. 
\end{lem}
\begin{proof}
Take any $t \in [\minsub{t}, \maxsub{t}]$. Then
\begin{equation*}
    \begin{split}
        \norm{f(t) - x_0} &\le \norm{f(t) - f(\minsub{t})} + \norm{f(\minsub{t}) - x_0} \\
        &\stackrel{(1)}{\le} \norm{t - \minsub{t}} + r \\ 
        &\le \norm{\maxsub{t} - \minsub{t}} + r \\
        &\le a - r + r = a
    \end{split}
\end{equation*}
where in (1) we used that $f \in \fM$. 
\end{proof}

\begin{lem}\label{lem:contFunInInt}
Let $f$ be Picard-Lindelöf with respect to $x_0, a, r, L, K$ and let $g \in \fM$. 
Then the function $[\minsub{t}, \maxsub{t}] \to E, t \mapsto f ((t, g(t)))$ is continuous. 
\end{lem}
\begin{proof}
Using the above lemma and the fact that $g$ is continuous on $s$, it suffices to show that $f$ is continuous on $[\minsub{t}, \maxsub{t}] \times (s \cap \closedBall{a}{x_0})$. 
This then follows from conditions \ref{PicLindI} and \ref{PicLindII} of $f$ being Picard-Lindelöf. 
\end{proof}

\begin{lem}
Let $f$ be Picard-Lindelöf with respect to $x_0, a, r, L, K$ and let $g \in \fM$. 
Then $C(f)$ is Lipschitz continuous with constant $L$. 
\end{lem}

\begin{proof}
Take $t_1, t_2 \in [\minsub{t}, \maxsub{t}]$. We assume that $t_2 \le t_1$
Then we get 
\begin{equation*}
    \begin{split}
        &\norm{C(f)(t_1) - C(f)(t_2)} \\
        =& \norm{x_0 + \int_{\minsub{t}}^{t_1} f(x, g(x)) \diff x - (x_0 + \int_{\minsub{t}}^{t_2} f(x, g(x)) \diff x)} \\ 
        \stackrel{(1)}{=}& \norm{\int_{t_2}^{t_1} f(x, g(x)) \diff x} \\ 
        \le& \int_{t_2}^{t_1} \norm{f(x, g(x))} \diff x \\ 
        \stackrel{(2)}{\le}& \int_{t_2}^{t_1} L \diff x \\ 
        =& L(t_1 - t_2)
    \end{split}
\end{equation*}
where in (1) we used that the function $x \mapsto f(x, g(x))$ is continuous by Lemma \ref{lem:contFunInInt} and thus integrable. 
In (2), we used Lemma \ref{lem:contFunInInt} and condition \ref{PicLindI} of $f$ being Picard-Lindelöf.
\end{proof}

\begin{lem}
Let $f$ be Picard-Lindelöf with respect to $x_0, a, r, L, K$ and let $g \in \fM$. 
Then $C(f)([\minsub{t}, \maxsub{t}]) \subseteq s$. 
\end{lem}

\begin{proof}
\todo[inline]{do this proof next}
\end{proof}